% ==================================================================
% id:       comb-013
% topic:    קומבינטוריקה
% subtopic: הסברים קומבינטוריים לזהויות
% solution: full
% ==================================================================
יהיו $n,k$ מספרים טבעיים, $1\le k\le n-1$. בשאלה נוכיח את \textbf{זהות פסקל}:
\[
\binom nk=\binom{n-1}{k-1}+\binom{n-1}k.
\]
\begin{enumerate}
\item הסבירו את הזהות במילים, בהסבר קומבינטורי.
\hint{מתוך $n$ תלמידים, שאחד מהם הוא דני, בוחרים $k$. מפרידים לשני מקרים: דני נבחר, או דני לא נבחר.}
\begin{solbox}
בכיתה $n$ תלמידים, ואחד מהם הוא דני. בוחרים $k$ תלמידים, בלי חשיבות לסדר ובלי חזרות. מצד אחד, יש $\binom nk$ דרכים לעשות זאת.

מצד שני, נחלק את הבחירות לשני מקרים זרים:
\begin{itemize}
\item \textbf{דני נבחר.} נשאר לבחור עוד $k-1$ תלמידים מתוך $n-1$ התלמידים האחרים: $\binom{n-1}{k-1}$ דרכים.
\item \textbf{דני לא נבחר.} בוחרים את כל $k$ התלמידים מתוך $n-1$ התלמידים האחרים: $\binom{n-1}k$ דרכים.
\end{itemize}
כל בחירה שייכת למקרה אחד בדיוק, ולכן לפי עקרון החיבור $\ans{\binom nk=\binom{n-1}{k-1}+\binom{n-1}k}$.
\end{solbox}
\item הוכיחו את הזהות אלגברית, בעזרת ההגדרה $\binom nk=\frac{n!}{k!\,(n-k)!}$.
\begin{solbox}
נכתוב את שני המחוברים באגף הימני לפי ההגדרה:
\[
\binom{n-1}{k-1}+\binom{n-1}k=\frac{(n-1)!}{(k-1)!\,(n-k)!}+\frac{(n-1)!}{k!\,(n-1-k)!}.
\]
נביא למכנה משותף $k!\,(n-k)!$: במחובר הראשון נכפיל מונה ומכנה ב-$k$ (כי $k\cdot(k-1)!=k!$), ובשני ב-$n-k$ (כי $(n-k)\cdot(n-k-1)!=(n-k)!$):
\[
\frac{k\cdot(n-1)!}{k!\,(n-k)!}+\frac{(n-k)\cdot(n-1)!}{k!\,(n-k)!}
=\frac{(n-1)!\,\big(k+(n-k)\big)}{k!\,(n-k)!}
=\frac{n\cdot(n-1)!}{k!\,(n-k)!}
=\frac{n!}{k!\,(n-k)!}=\binom nk.
\]
\end{solbox}
\end{enumerate}
